\documentclass[12pt,a4paper]{article}

% ==================== GÓI CƠ BẢN ====================
\usepackage[utf8]{inputenc}
\usepackage[T5]{fontenc}
\usepackage[vietnamese]{babel}
\usepackage{geometry}
\usepackage{parskip}
\usepackage{enumitem}
\usepackage{longtable}
\usepackage{array}
\usepackage{xcolor}
\usepackage{hyperref}
\usepackage{fancyhdr}
\usepackage{lastpage}

% ==================== ĐỊNH DẠNG TRANG ====================
\geometry{
    top=2cm,
    bottom=2cm,
    left=2.5cm,
    right=2cm
}

\setlength{\parindent}{0pt}
\setlength{\parskip}{0.7em}
\setlist[itemize]{leftmargin=1.8em, itemsep=0.35em}
\setlist[enumerate]{leftmargin=1.8em, itemsep=0.35em}

\definecolor{ailmsblue}{RGB}{25,84,166}
\definecolor{ailmsgray}{RGB}{90,90,90}
\definecolor{ailmsred}{RGB}{180,40,40}

\hypersetup{
    colorlinks=true,
    linkcolor=ailmsblue,
    urlcolor=ailmsblue,
    citecolor=ailmsblue,
    pdftitle={Chính sách Khóa học và Đào tạo AILMS},
    pdfauthor={Ban Quản lý Đào tạo AILMS}
}

% ==================== HEADER / FOOTER ====================
\pagestyle{fancy}
\fancyhf{}
\fancyhead[L]{\textcolor{ailmsgray}{AILMS}}
\fancyhead[R]{\textcolor{ailmsgray}{Chính sách Khóa học \& Đào tạo}}
\fancyfoot[C]{Trang \thepage/\pageref{LastPage}}
\renewcommand{\headrulewidth}{0.4pt}
\renewcommand{\footrulewidth}{0.4pt}

% ==================== THÔNG TIN TÀI LIỆU ====================
\title{
    \textbf{\LARGE CHÍNH SÁCH KHÓA HỌC\\
    VÀ HOẠT ĐỘNG ĐÀO TẠO AILMS}
}
\author{Ban Quản lý Đào tạo AILMS}
\date{Phiên bản 1.0 -- Ngày hiệu lực: \today}

\begin{document}

\maketitle

\begin{center}
    \textit{Tài liệu này quy định quyền lợi, nghĩa vụ và nguyên tắc tổ chức
    hoạt động đào tạo đối với các khóa học được cung cấp trên hệ thống AILMS.}
\end{center}

\vspace{0.8cm}

\begin{longtable}{|>{\bfseries}p{4cm}|p{9.5cm}|}
    \hline
    Tên tài liệu &
    Chính sách Khóa học và Hoạt động đào tạo AILMS \\
    \hline
    Phiên bản & 1.0 \\
    \hline
    Đơn vị quản lý & Ban Quản lý Đào tạo AILMS \\
    \hline
    Đối tượng áp dụng &
    Học viên, giảng viên, trợ giảng, người giám hộ,
    quản trị viên đào tạo và các bên tham gia tổ chức khóa học \\
    \hline
    Phạm vi áp dụng &
    Khóa học tự học, lớp học nhóm, đào tạo một-một,
    bài kiểm tra, bài tập và chứng chỉ trên AILMS \\
    \hline
    Ngày hiệu lực & \today \\
    \hline
\end{longtable}

\tableofcontents
\newpage

% =========================================================
\section{Mục đích và phạm vi áp dụng}

Chính sách này được ban hành nhằm:

\begin{itemize}
    \item Quy định điều kiện đăng ký và sử dụng khóa học.
    \item Xác định quyền và nghĩa vụ của học viên, giảng viên và trợ giảng.
    \item Thiết lập nguyên tắc tổ chức các hình thức đào tạo.
    \item Quy định việc quản lý lịch học, điểm danh, bài tập và đánh giá.
    \item Bảo đảm tính minh bạch về thời hạn, quyền lợi và điều kiện hoàn thành.
    \item Thiết lập cơ chế xử lý đổi lịch, hủy buổi, bảo lưu và khiếu nại.
\end{itemize}

Chính sách áp dụng đối với toàn bộ khóa học và gói đào tạo được công bố
trên AILMS, trừ trường hợp từng khóa học có quy định riêng được hiển thị rõ
trước khi học viên đăng ký.

% =========================================================
\section{Giải thích thuật ngữ}

Trong tài liệu này, các thuật ngữ được hiểu như sau:

\begin{itemize}
    \item \textbf{Khóa học:} Tập hợp nội dung đào tạo gồm bài học,
    tài liệu, bài tập, bài kiểm tra và các hoạt động liên quan.

    \item \textbf{Gói đào tạo:} Phương án cung cấp một khóa học với
    hình thức, mức giá, thời hạn và quyền lợi cụ thể.

    \item \textbf{Tự học -- Self-Paced:} Hình thức học viên chủ động
    truy cập nội dung và tự sắp xếp tiến độ.

    \item \textbf{Lớp học nhóm -- Group Class:} Hình thức nhiều học viên
    cùng tham gia một lớp có giảng viên, sĩ số và lịch học cụ thể.

    \item \textbf{Đào tạo một-một -- One-on-One:} Hình thức một học viên
    học hoặc tư vấn trực tiếp với một giảng viên theo lịch đã thống nhất.

    \item \textbf{Enrollment:} Bản ghi xác nhận quyền tham gia khóa học
    hoặc gói đào tạo của học viên.

    \item \textbf{Waitlist:} Danh sách chờ dành cho học viên chưa thể vào lớp
    do lớp đã đủ sĩ số hoặc chưa đáp ứng điều kiện tiếp nhận.

    \item \textbf{Buổi học:} Một phiên giảng dạy có thời gian bắt đầu,
    thời lượng, hình thức và người phụ trách xác định.

    \item \textbf{Tiến độ học tập:} Tỷ lệ hoặc trạng thái hoàn thành bài học,
    nội dung, bài tập và bài kiểm tra của học viên.

    \item \textbf{Pass Score:} Điểm tối thiểu cần đạt để vượt qua
    một bài kiểm tra hoặc hoàn thành điều kiện của khóa học.
\end{itemize}

% =========================================================
\section{Thông tin khóa học và gói đào tạo}

Trước khi đăng ký hoặc thanh toán, học viên phải được cung cấp các thông tin chính:

\begin{itemize}
    \item Tên và mô tả khóa học.
    \item Mục tiêu và kết quả học tập dự kiến.
    \item Đối tượng phù hợp và yêu cầu đầu vào.
    \item Trình độ khóa học.
    \item Danh mục và giảng viên tạo khóa học.
    \item Hình thức đào tạo.
    \item Nội dung hoặc chương trình học.
    \item Mức giá và quyền lợi của từng gói.
    \item Thời hạn truy cập.
    \item Điều kiện hoàn thành và cấp chứng chỉ.
    \item Điều kiện hoàn tiền hoặc hủy đăng ký.
\end{itemize}

Nếu thông tin giữa trang giới thiệu và trang thanh toán có khác biệt,
thông tin được hiển thị tại thời điểm học viên xác nhận đơn hàng được ưu tiên,
trừ trường hợp có lỗi hiển thị rõ ràng.

% =========================================================
\section{Điều kiện đăng ký khóa học}

\subsection{Điều kiện chung}

Học viên được đăng ký khóa học khi:

\begin{itemize}
    \item Có tài khoản hợp lệ và không bị đình chỉ.
    \item Đáp ứng yêu cầu đầu vào nếu khóa học có quy định.
    \item Đồng ý với chính sách khóa học và điều khoản liên quan.
    \item Hoàn tất thanh toán hoặc được cấp quyền theo chương trình hợp lệ.
    \item Cung cấp đủ thông tin cần thiết đối với gói lớp nhóm hoặc một-một.
\end{itemize}

\subsection{Thời điểm phát sinh quyền học}

Việc tạo đơn hàng hoặc chuyển đến cổng thanh toán chưa làm phát sinh quyền học.

Quyền học được cấp sau khi:

\begin{enumerate}
    \item Giao dịch được Backend xác minh thành công; hoặc
    \item Đơn hàng miễn phí được xác nhận hợp lệ; hoặc
    \item Người có thẩm quyền cấp quyền theo chương trình đào tạo.
\end{enumerate}

\subsection{Đăng ký trùng lặp}

Hệ thống có thể từ chối đăng ký nếu học viên:

\begin{itemize}
    \item Đã sở hữu cùng gói đào tạo còn hiệu lực.
    \item Đang có đơn hàng chờ thanh toán cho cùng gói.
    \item Đã tham gia lớp nhóm được liên kết với gói đó.
    \item Không đủ điều kiện tham gia hoặc đã bị hạn chế quyền học.
\end{itemize}

% =========================================================
\section{Thời hạn và phạm vi truy cập}

Thời hạn truy cập được xác định theo từng gói đào tạo, chẳng hạn:

\begin{itemize}
    \item Truy cập trong một số ngày kể từ ngày kích hoạt.
    \item Truy cập đến ngày kết thúc khóa học hoặc lớp học.
    \item Truy cập dài hạn trong thời gian dịch vụ còn được duy trì.
\end{itemize}

Khái niệm ``trọn đời'' nếu được sử dụng được hiểu là thời gian nội dung
còn được AILMS cung cấp hợp pháp và hệ thống còn duy trì dịch vụ,
không đồng nghĩa với cam kết tồn tại vĩnh viễn.

Sau khi hết hạn:

\begin{itemize}
    \item Học viên có thể không còn quyền mở nội dung trả phí.
    \item Tiến độ và kết quả trước đó có thể tiếp tục được lưu.
    \item Học viên có thể gia hạn nếu gói đào tạo hỗ trợ.
    \item Chứng chỉ đã cấp hợp lệ không mặc nhiên bị hủy do gói hết hạn.
\end{itemize}

% =========================================================
\section{Hình thức tự học}

\subsection{Quyền lợi}

Học viên tự học có thể được:

\begin{itemize}
    \item Truy cập video, bài đọc và tài liệu thuộc gói đã mua.
    \item Thực hiện bài tập và bài kiểm tra trực tuyến.
    \item Theo dõi tiến độ học tập.
    \item Sử dụng AI Tutor trong giới hạn được áp dụng.
    \item Nhận chứng chỉ nếu hoàn thành đủ điều kiện.
\end{itemize}

\subsection{Trách nhiệm của học viên}

Học viên có trách nhiệm:

\begin{itemize}
    \item Chủ động xây dựng kế hoạch học tập.
    \item Hoàn thành nội dung trong thời hạn truy cập.
    \item Không chia sẻ tài khoản hoặc nội dung khóa học.
    \item Tự kiểm tra thiết bị và kết nối mạng.
    \item Tuân thủ quy định về bài tập, kiểm tra và bản quyền.
\end{itemize}

\subsection{Ghi nhận tiến độ}

Hệ thống có thể ghi nhận tiến độ dựa trên:

\begin{itemize}
    \item Trạng thái thái mở và hoàn thành bài học.
    \item Thời lượng xem nội dung.
    \item Kết quả bài tập hoặc bài kiểm tra.
    \item Các yêu cầu hoàn thành riêng của từng bài học.
\end{itemize}

Việc mở một bài học không mặc nhiên được xem là đã hoàn thành
nếu học viên chưa đáp ứng điều kiện theo cấu hình khóa học.

% =========================================================
\section{Hình thức lớp học nhóm}

\subsection{Thông tin lớp học}

Mỗi lớp nhóm cần xác định:

\begin{itemize}
    \item Khóa học và gói đào tạo liên quan.
    \item Giảng viên hoặc trợ giảng phụ trách.
    \item Sĩ số tối đa.
    \item Thời gian bắt đầu và kết thúc.
    \item Lịch học dự kiến.
    \item Hình thức học trực tuyến hoặc trực tiếp.
    \item Nội quy và điều kiện hoàn thành.
\end{itemize}

\subsection{Xếp học viên vào lớp}

Sau khi thanh toán thành công:

\begin{itemize}
    \item Học viên được thêm vào lớp nếu lớp còn chỗ và đủ điều kiện.
    \item Hệ thống gửi thông báo cho học viên và giảng viên liên quan.
    \item Quyền truy cập không gian lớp học được cấp theo vai trò học viên.
    \item Học viên không được truy cập chức năng quản trị dành cho giảng viên.
\end{itemize}

Nếu lớp đã đủ sĩ số, học viên có thể được:

\begin{itemize}
    \item Đưa vào Waitlist.
    \item Chuyển sang lớp tương đương.
    \item Được đề xuất lịch hoặc gói khác.
    \item Yêu cầu hoàn tiền nếu hệ thống không thể cung cấp quyền lợi phù hợp.
\end{itemize}

\subsection{Danh sách chờ}

Việc đưa vào Waitlist không đồng nghĩa học viên đã chắc chắn có chỗ.

Khi có chỗ trống, hệ thống có thể xét theo:

\begin{itemize}
    \item Thứ tự tham gia danh sách chờ.
    \item Điều kiện đầu vào.
    \item Khả năng tham gia lịch học.
    \item Trạng thái thanh toán và tài khoản.
\end{itemize}

Học viên phải xác nhận tham gia trong thời hạn được thông báo.
Nếu không xác nhận, quyền ưu tiên có thể được chuyển cho người tiếp theo.

\subsection{Thay đổi sĩ số}

Sĩ số tối đa được xác định theo năng lực tổ chức của từng lớp.
AILMS không tự ý vượt quá sĩ số đã công bố nếu việc đó ảnh hưởng đáng kể
đến chất lượng hoặc quyền lợi học viên.

% =========================================================
\section{Hình thức đào tạo một-một}

\subsection{Yêu cầu đăng ký}

Khi đăng ký gói một-một, học viên có thể được yêu cầu cung cấp:

\begin{itemize}
    \item Mục tiêu học tập.
    \item Kiến thức và trình độ hiện tại.
    \item Nội dung đang gặp khó khăn.
    \item Kết quả hoặc chứng chỉ mong muốn.
    \item Khung giờ có thể tham gia.
    \item Yêu cầu hợp lý đối với giảng viên.
\end{itemize}

Thông tin này chỉ được sử dụng để tư vấn, xếp lịch và tổ chức đào tạo.

\subsection{Ghép giảng viên}

Hệ thống hoặc bộ phận phụ trách có thể đề xuất giảng viên dựa trên:

\begin{itemize}
    \item Chuyên môn và khóa học phụ trách.
    \item Lịch rảnh của giảng viên.
    \item Khung giờ của học viên.
    \item Trình độ và mục tiêu học tập.
    \item Số buổi còn khả dụng.
\end{itemize}

Việc mua gói không mặc nhiên xác nhận một giảng viên hoặc khung giờ cụ thể
cho đến khi các bên thống nhất.

\subsection{Tổ chức lớp một-một}

Sau khi học viên và giảng viên trao đổi, giảng viên hoặc trợ giảng có quyền phù hợp
có thể tạo lớp, tạo lịch và quản lý các buổi học một-một.

Bộ phận nhân sự hoặc quản trị chỉ can thiệp khi:

\begin{itemize}
    \item Không tìm được giảng viên phù hợp.
    \item Giảng viên không thể tiếp tục đảm nhiệm.
    \item Có tranh chấp hoặc yêu cầu thay đổi đặc biệt.
    \item Cần kiểm tra điều kiện hợp đồng hoặc phân công chuyên môn.
\end{itemize}

Mô hình này nhằm hạn chế việc mọi thao tác tổ chức lớp đều phụ thuộc vào nhân sự.

% =========================================================
\section{Lịch học và thay đổi lịch}

\subsection{Công bố lịch học}

Lịch học phải thể hiện tối thiểu:

\begin{itemize}
    \item Ngày và giờ bắt đầu.
    \item Thời lượng dự kiến.
    \item Giảng viên phụ trách.
    \item Hình thức và địa điểm hoặc liên kết tham gia.
    \item Trạng thái của buổi học.
\end{itemize}

Thời gian trên hệ thống được hiển thị theo múi giờ được cấu hình
hoặc múi giờ được thông báo cho người dùng.

\subsection{Yêu cầu đổi lịch của học viên}

Đối với buổi học một-một, học viên nên gửi yêu cầu đổi lịch
ít nhất 24 giờ trước giờ bắt đầu, trừ khi gói đào tạo quy định khác.

Yêu cầu đổi lịch:

\begin{itemize}
    \item Phải nêu lý do và đề xuất khung giờ thay thế.
    \item Chỉ có hiệu lực khi được giảng viên hoặc hệ thống xác nhận.
    \item Có thể bị giới hạn số lần trong một chu kỳ hoặc gói học.
    \item Không bảo đảm giữ nguyên giảng viên nếu lịch mới không phù hợp.
\end{itemize}

\subsection{Thay đổi lịch của giảng viên}

Nếu giảng viên cần đổi lịch, giảng viên phải thông báo sớm
và đề xuất phương án thay thế.

Học viên có quyền:

\begin{itemize}
    \item Chấp nhận lịch mới.
    \item Đề xuất khung giờ khác.
    \item Yêu cầu giảng viên thay thế nếu việc đổi lịch lặp lại.
    \item Yêu cầu xem xét hoàn tiền phần chưa sử dụng khi hệ thống
    không thể tiếp tục cung cấp dịch vụ.
\end{itemize}

% =========================================================
\section{Vắng mặt và hủy buổi học}

\subsection{Học viên vắng mặt}

Buổi học có thể được tính là đã sử dụng nếu:

\begin{itemize}
    \item Học viên không tham gia mà không thông báo.
    \item Học viên báo hủy sau thời hạn cho phép.
    \item Học viên đến quá muộn khiến buổi học không thể thực hiện hợp lý.
\end{itemize}

Trường hợp bất khả kháng hoặc sự cố có bằng chứng phù hợp
có thể được xem xét riêng.

\subsection{Giảng viên vắng mặt}

Nếu giảng viên vắng mặt hoặc hủy buổi không hợp lệ:

\begin{itemize}
    \item Buổi học không được trừ vào quyền lợi của học viên.
    \item Hệ thống phải tổ chức lịch học bù hoặc giảng viên thay thế.
    \item Học viên được thông báo về phương án xử lý.
    \item Yêu cầu hoàn tiền có thể được xem xét nếu sự việc tái diễn.
\end{itemize}

\subsection{Sự cố kỹ thuật}

Nếu buổi học bị gián đoạn do lỗi hệ thống hoặc hạ tầng thuộc trách nhiệm AILMS,
buổi học có thể được kéo dài, tổ chức lại hoặc bù quyền lợi.

Nếu sự cố đến từ thiết bị hoặc kết nối của một bên,
các bên phối hợp ghi nhận tình trạng và thống nhất phương án phù hợp.

% =========================================================
\section{Điểm danh}

Điểm danh có thể được thực hiện bằng:

\begin{itemize}
    \item Xác nhận thủ công của giảng viên hoặc trợ giảng.
    \item Ghi nhận thời gian tham gia buổi học trực tuyến.
    \item Mã điểm danh hoặc cơ chế xác thực phù hợp.
    \item Kết hợp nhiều nguồn thông tin để xử lý tranh chấp.
\end{itemize}

Các trạng thái điểm danh có thể bao gồm:

\begin{itemize}
    \item Có mặt.
    \item Đi muộn.
    \item Vắng có phép.
    \item Vắng không phép.
    \item Chưa xác nhận.
\end{itemize}

Học viên có quyền phản hồi nếu cho rằng trạng thái điểm danh không chính xác.

% =========================================================
\section{Không gian lớp học và tương tác}

Thành viên lớp học có thể được phép:

\begin{itemize}
    \item Xem thông báo và lịch học.
    \item Đăng câu hỏi hoặc bài thảo luận.
    \item Bình luận trong nội dung được cho phép.
    \item Xem và nộp bài tập.
    \item Truy cập tài liệu được chia sẻ cho lớp.
\end{itemize}

Học viên không được:

\begin{itemize}
    \item Đăng nội dung spam hoặc không liên quan.
    \item Chia sẻ mã độc, liên kết lừa đảo hoặc tệp nguy hiểm.
    \item Công khai thông tin cá nhân của thành viên khác.
    \item Xúc phạm, quấy rối hoặc gây gián đoạn hoạt động lớp.
    \item Đăng đáp án bài kiểm tra khi chưa được cho phép.
    \item Quảng cáo hoặc bán hàng trong không gian lớp học.
\end{itemize}

Giảng viên hoặc trợ giảng có thể ẩn, khóa bình luận hoặc báo cáo nội dung vi phạm.
Các thao tác kiểm duyệt quan trọng nên được ghi nhận để phục vụ truy vết.

% =========================================================
\section{Tài liệu và nội dung bài học}

Tài liệu khóa học chỉ được sử dụng cho mục đích học tập cá nhân
trong phạm vi quyền truy cập được cấp.

Học viên không được:

\begin{itemize}
    \item Sao chép hoặc phát tán tài liệu trái phép.
    \item Ghi hình hoặc phát trực tiếp buổi học khi chưa được đồng ý.
    \item Chia sẻ liên kết truy cập riêng tư.
    \item Tải xuống nội dung bằng công cụ vượt qua giới hạn kỹ thuật.
    \item Sử dụng tài liệu làm dữ liệu huấn luyện AI bên ngoài khi chưa được phép.
\end{itemize}

Tài liệu có thể được cập nhật để sửa lỗi, cải thiện chất lượng
hoặc đáp ứng thay đổi của chương trình đào tạo.

% =========================================================
\section{Bài tập}

\subsection{Các loại bài tập}

Khóa học có thể sử dụng:

\begin{itemize}
    \item Bài tập tự luận.
    \item Bài tập lập trình.
    \item Bài tập tải tệp.
    \item Bài tập nhóm.
    \item Bài tập thực hành hoặc dự án.
\end{itemize}

\subsection{Nộp bài}

Học viên có trách nhiệm:

\begin{itemize}
    \item Nộp đúng định dạng và thời hạn.
    \item Kiểm tra tệp trước khi xác nhận nộp.
    \item Bảo đảm bài làm là sản phẩm hợp lệ của mình.
    \item Lưu bản sao bài làm khi cần thiết.
\end{itemize}

Thời điểm nộp được xác định theo thời gian hệ thống ghi nhận thành công,
không phải thời điểm học viên bắt đầu tải tệp lên.

\subsection{Nộp muộn}

Tùy cấu hình khóa học, bài nộp muộn có thể:

\begin{itemize}
    \item Bị từ chối.
    \item Bị trừ điểm.
    \item Được chấp nhận sau khi giảng viên phê duyệt.
    \item Được gia hạn trong trường hợp có lý do hợp lệ.
\end{itemize}

% =========================================================
\section{Bài kiểm tra và đánh giá}

\subsection{Hình thức đánh giá}

Khóa học có thể sử dụng:

\begin{itemize}
    \item Câu hỏi một lựa chọn.
    \item Câu hỏi nhiều lựa chọn.
    \item Câu hỏi đúng hoặc sai.
    \item Câu hỏi điền khuyết.
    \item Câu hỏi tự luận.
    \item Bài tập hoặc dự án thực hành.
\end{itemize}

\subsection{Chấm điểm}

Các câu hỏi có đáp án xác định có thể được chấm tự động.

Bài tự luận, bài tập, dự án và các nội dung yêu cầu đánh giá chuyên môn
có thể do giảng viên hoặc trợ giảng chấm thủ công.

AI có thể hỗ trợ:

\begin{itemize}
    \item Tạo bản nháp câu hỏi.
    \item Gợi ý tiêu chí đánh giá.
    \item Phân tích sơ bộ bài làm.
    \item Đề xuất phản hồi cho người chấm.
\end{itemize}

Kết quả có ảnh hưởng quan trọng không nên chỉ dựa hoàn toàn vào AI
nếu chưa có bước kiểm tra phù hợp của con người.

\subsection{Phúc khảo}

Học viên có thể yêu cầu xem xét lại kết quả khi:

\begin{itemize}
    \item Điểm không khớp với bài đã nộp.
    \item Hệ thống chấm tự động có dấu hiệu sai.
    \item Có sai sót trong tiêu chí hoặc quá trình chấm.
\end{itemize}

Yêu cầu phúc khảo phải được gửi trong thời hạn do khóa học công bố
và kèm theo lý do cụ thể.

% =========================================================
\section{Điều kiện hoàn thành khóa học}

Tùy từng khóa học, học viên có thể phải đáp ứng một hoặc nhiều điều kiện:

\begin{itemize}
    \item Hoàn thành tỷ lệ nội dung tối thiểu.
    \item Đạt Pass Score của bài kiểm tra.
    \item Nộp đủ bài tập bắt buộc.
    \item Đạt yêu cầu điểm danh.
    \item Hoàn thành dự án cuối khóa.
    \item Không có vi phạm học thuật chưa được xử lý.
\end{itemize}

Điều kiện hoàn thành phải được công bố trong thông tin khóa học
hoặc không gian học tập.

% =========================================================
\section{Chứng chỉ}

\subsection{Điều kiện cấp}

Chứng chỉ có thể được cấp khi:

\begin{itemize}
    \item Học viên hoàn thành đủ điều kiện.
    \item Kết quả đã được xác nhận.
    \item Tài khoản và thông tin học viên hợp lệ.
    \item Không có vi phạm nghiêm trọng ảnh hưởng đến kết quả.
\end{itemize}

\subsection{Thông tin chứng chỉ}

Chứng chỉ có thể bao gồm:

\begin{itemize}
    \item Họ tên học viên.
    \item Tên khóa học.
    \item Ngày hoàn thành.
    \item Mã chứng chỉ.
    \item Đơn vị cấp.
    \item Thông tin xác minh.
\end{itemize}

\subsection{Giá trị của chứng chỉ}

Chứng chỉ AILMS xác nhận học viên đã hoàn thành điều kiện của khóa học.

Chứng chỉ không mặc nhiên tương đương với:

\begin{itemize}
    \item Văn bằng quốc gia.
    \item Chứng chỉ hành nghề.
    \item Cam kết tuyển dụng.
    \item Bảo đảm năng lực trong mọi tình huống thực tế.
\end{itemize}

trừ khi khóa học có công bố và căn cứ pháp lý riêng.

\subsection{Thu hồi chứng chỉ}

Chứng chỉ có thể bị thu hồi nếu:

\begin{itemize}
    \item Phát hiện gian lận nghiêm trọng.
    \item Thông tin nhận dạng bị giả mạo.
    \item Chứng chỉ được cấp do lỗi hệ thống hoặc dữ liệu sai.
    \item Có quyết định hợp lệ từ đơn vị có thẩm quyền.
\end{itemize}

Việc thu hồi phải có lý do và được ghi nhận trên hệ thống.

% =========================================================
\section{Liêm chính học thuật và chống gian lận}

Các hành vi gian lận bao gồm:

\begin{itemize}
    \item Nhờ người khác học, thi hoặc nộp bài thay.
    \item Sử dụng nhiều tài khoản để thao túng kết quả.
    \item Sao chép bài làm mà không trích dẫn.
    \item Mua bán đáp án hoặc bài tập.
    \item Chia sẻ câu hỏi kiểm tra trái phép.
    \item Can thiệp vào dữ liệu điểm hoặc tiến độ.
    \item Dùng công cụ tự động để giả lập hoạt động học.
    \item Sử dụng AI trái với quy định của bài tập.
\end{itemize}

Hệ thống có thể sử dụng thông tin kỹ thuật để phát hiện dấu hiệu bất thường,
nhưng không nên kết luận vi phạm chỉ dựa trên một tín hiệu tự động duy nhất.

Học viên phải được tạo cơ hội giải trình trước khi áp dụng hình thức xử lý nghiêm trọng,
trừ trường hợp cần khóa tạm thời để bảo vệ hệ thống hoặc bằng chứng.

% =========================================================
\section{Sử dụng AI trong học tập}

AI Tutor có thể hỗ trợ:

\begin{itemize}
    \item Giải thích bài học.
    \item Gợi ý tài liệu.
    \item Đề xuất lộ trình.
    \item Hỗ trợ ôn tập.
    \item Phản hồi sơ bộ đối với câu trả lời.
\end{itemize}

Học viên không được sử dụng AI để:

\begin{itemize}
    \item Tạo toàn bộ bài làm rồi khai nhận là sản phẩm cá nhân,
    nếu bài tập không cho phép.
    \item Gian lận trong bài kiểm tra.
    \item Truy xuất đáp án hoặc dữ liệu chưa được phép.
    \item Làm thay hoạt động đánh giá năng lực cá nhân.
\end{itemize}

Giảng viên nên công bố rõ đối với từng bài tập:

\begin{itemize}
    \item Có được sử dụng AI hay không.
    \item Được sử dụng AI ở mức độ nào.
    \item Có phải khai báo hoặc trích dẫn việc sử dụng AI hay không.
\end{itemize}

% =========================================================
\section{Bảo lưu và chuyển lớp}

\subsection{Bảo lưu}

Học viên có thể yêu cầu bảo lưu trong các trường hợp:

\begin{itemize}
    \item Sức khỏe không bảo đảm.
    \item Có sự kiện bất khả kháng.
    \item Không thể tiếp tục tham gia lịch học trong một khoảng thời gian.
    \item Có lý do khác được đơn vị đào tạo chấp thuận.
\end{itemize}

Thời gian và số lần bảo lưu phụ thuộc vào từng gói đào tạo.

\subsection{Chuyển lớp}

Học viên có thể được xem xét chuyển lớp khi:

\begin{itemize}
    \item Lịch học không còn phù hợp.
    \item Trình độ lớp không phù hợp sau đánh giá.
    \item Giảng viên hoặc lớp hiện tại không thể tiếp tục.
    \item Lớp khác còn chỗ và đáp ứng điều kiện.
\end{itemize}

Nếu lớp mới có mức giá hoặc quyền lợi khác, hệ thống có thể yêu cầu
xử lý phần chênh lệch theo chính sách tại thời điểm chuyển.

% =========================================================
\section{Tạm dừng, hủy hoặc thay đổi khóa học}

AILMS có thể tạm dừng hoặc điều chỉnh khóa học nếu:

\begin{itemize}
    \item Không đủ số lượng học viên tối thiểu.
    \item Giảng viên không thể tiếp tục giảng dạy.
    \item Nội dung cần cập nhật hoặc có vấn đề bản quyền.
    \item Xảy ra sự cố kỹ thuật hoặc sự kiện bất khả kháng.
    \item Khóa học không còn đáp ứng tiêu chuẩn chất lượng.
\end{itemize}

Trong trường hợp quyền lợi học viên bị ảnh hưởng đáng kể,
AILMS có thể cung cấp một trong các phương án:

\begin{itemize}
    \item Chuyển sang lớp hoặc khóa học tương đương.
    \item Thay giảng viên.
    \item Gia hạn thời gian truy cập.
    \item Bổ sung buổi học hoặc quyền lợi.
    \item Hoàn tiền cho phần dịch vụ chưa được cung cấp.
\end{itemize}

% =========================================================
\section{Chính sách hoàn tiền liên quan đến đào tạo}

Điều kiện hoàn tiền cụ thể được áp dụng theo chính sách hoàn tiền chung
và thông tin của từng gói.

\subsection{Gói tự học}

Yêu cầu có thể được xem xét dựa trên:

\begin{itemize}
    \item Thời gian kể từ ngày kích hoạt.
    \item Tỷ lệ nội dung đã hoàn thành.
    \item Tài liệu đã tải xuống.
    \item Bài kiểm tra hoặc chứng chỉ đã được thực hiện.
\end{itemize}

\subsection{Gói lớp học nhóm}

Yêu cầu có thể được xem xét khi:

\begin{itemize}
    \item Lớp bị hủy.
    \item Hệ thống không thể xếp lớp phù hợp.
    \item Lịch học thay đổi đáng kể và không có phương án thay thế.
    \item Phần lớn quyền lợi của gói chưa được sử dụng.
\end{itemize}

\subsection{Gói một-một}

Yêu cầu có thể được xem xét đối với số buổi chưa sử dụng khi:

\begin{itemize}
    \item Hệ thống không thể cung cấp giảng viên phù hợp.
    \item Giảng viên liên tục hủy lịch và không có người thay thế.
    \item Việc tổ chức đào tạo không thể tiếp tục do lỗi thuộc phía cung cấp dịch vụ.
\end{itemize}

Các buổi đã thực hiện hoặc bị tính là đã sử dụng hợp lệ
thường không thuộc phạm vi hoàn tiền.

% =========================================================
\section{Trách nhiệm của học viên}

Học viên có trách nhiệm:

\begin{itemize}
    \item Cung cấp thông tin chính xác.
    \item Chủ động theo dõi lịch và thông báo.
    \item Tham gia đúng giờ.
    \item Hoàn thành bài tập và đánh giá đúng hạn.
    \item Tôn trọng giảng viên và thành viên lớp.
    \item Bảo vệ tài khoản và tài liệu.
    \item Báo cáo sớm khi phát hiện sai sót.
\end{itemize}

% =========================================================
\section{Trách nhiệm của giảng viên}

Giảng viên có trách nhiệm:

\begin{itemize}
    \item Cung cấp nội dung đúng phạm vi và mục tiêu khóa học.
    \item Công bố yêu cầu học tập rõ ràng.
    \item Thực hiện lịch giảng dạy đã xác nhận.
    \item Chấm và phản hồi trong thời gian hợp lý.
    \item Tôn trọng quyền riêng tư của học viên.
    \item Không tiết lộ điểm, bài làm hoặc thông tin cá nhân trái phép.
    \item Báo cáo các vấn đề ảnh hưởng đến chất lượng khóa học.
\end{itemize}

% =========================================================
\section{Trách nhiệm của trợ giảng}

Trợ giảng có thể:

\begin{itemize}
    \item Hỗ trợ quản lý lớp và lịch học.
    \item Giải đáp câu hỏi trong phạm vi được giao.
    \item Hỗ trợ chấm bài theo tiêu chí của giảng viên.
    \item Kiểm duyệt thảo luận lớp.
    \item Ghi nhận điểm danh và sự cố.
\end{itemize}

Trợ giảng không được thực hiện quyền của giảng viên hoặc quản trị viên
nếu chưa được phân quyền hợp lệ.

% =========================================================
\section{Đánh giá chất lượng khóa học}

Học viên có thể đánh giá khóa học sau khi đạt điều kiện tham gia hợp lệ.

Nội dung đánh giá phải:

\begin{itemize}
    \item Phản ánh trải nghiệm thực tế.
    \item Không chứa thông tin sai lệch có chủ ý.
    \item Không quấy rối hoặc xúc phạm cá nhân.
    \item Không chứa quảng cáo hoặc nội dung không liên quan.
\end{itemize}

AILMS không được tự ý thay đổi nội dung đánh giá hợp lệ
chỉ vì đánh giá mang tính tiêu cực.

Đánh giá có thể bị ẩn nếu vi phạm quy tắc cộng đồng,
có dấu hiệu gian lận hoặc không xuất phát từ trải nghiệm học tập hợp lệ.

% =========================================================
\section{Khiếu nại và giải quyết tranh chấp đào tạo}

Học viên có thể khiếu nại về:

\begin{itemize}
    \item Tiến độ hoặc quyền truy cập.
    \item Điểm danh.
    \item Kết quả bài tập hoặc kiểm tra.
    \item Lịch học và việc vắng mặt.
    \item Chất lượng giảng dạy.
    \item Chứng chỉ.
    \item Hoàn tiền hoặc quyền lợi gói đào tạo.
\end{itemize}

Yêu cầu nên kèm theo:

\begin{itemize}
    \item Mã học viên.
    \item Mã khóa học, lớp học hoặc đơn hàng.
    \item Mô tả cụ thể sự việc.
    \item Thời điểm phát sinh.
    \item Bằng chứng liên quan nếu có.
\end{itemize}

AILMS hướng tới xác nhận tiếp nhận yêu cầu trong vòng 48 giờ làm việc.
Thời gian xử lý cuối cùng phụ thuộc vào tính chất và mức độ phức tạp của vụ việc.

% =========================================================
\section{Sự kiện bất khả kháng}

Các sự kiện bất khả kháng có thể bao gồm:

\begin{itemize}
    \item Thiên tai, dịch bệnh hoặc hỏa hoạn.
    \item Mất điện hoặc gián đoạn Internet trên diện rộng.
    \item Sự cố nghiêm trọng của nền tảng hội nghị trực tuyến.
    \item Quyết định của cơ quan có thẩm quyền.
    \item Tình huống khác nằm ngoài khả năng kiểm soát hợp lý.
\end{itemize}

Các bên có trách nhiệm thông báo và phối hợp lựa chọn phương án
học bù, đổi lịch, bảo lưu hoặc xử lý quyền lợi phù hợp.

% =========================================================
\section{Sửa đổi chính sách}

AILMS có thể cập nhật chính sách này khi:

\begin{itemize}
    \item Có thay đổi về hình thức đào tạo.
    \item Có thay đổi về chức năng hệ thống.
    \item Cần điều chỉnh quy trình học tập hoặc đánh giá.
    \item Có yêu cầu từ pháp luật hoặc cơ quan quản lý.
\end{itemize}

Phiên bản mới phải ghi rõ ngày hiệu lực.

Nếu thay đổi ảnh hưởng đáng kể đến gói mà học viên đã mua,
AILMS sẽ thực hiện biện pháp thông báo và xử lý quyền lợi phù hợp.

% =========================================================
\section{Thông tin liên hệ}

Các yêu cầu liên quan đến đào tạo có thể được gửi tới:

\begin{itemize}
    \item \textbf{Đơn vị phụ trách:} Ban Quản lý Đào tạo AILMS.
    \item \textbf{Email hỗ trợ:}
    \href{mailto:training@ailms.example}{training@ailms.example}.
    \item \textbf{Kênh hỗ trợ:} Chức năng gửi yêu cầu trên hệ thống.
    \item \textbf{Thời gian hỗ trợ:} Theo lịch làm việc được công bố.
\end{itemize}

\textcolor{ailmsred}{
\textit{Lưu ý: Thay địa chỉ thư điện tử mẫu bằng thông tin chính thức
trước khi công bố tài liệu.}
}

% =========================================================
\section{Hiệu lực thi hành}

Chính sách này có hiệu lực kể từ ngày được công bố.

Học viên, giảng viên, trợ giảng và các đơn vị liên quan có trách nhiệm
đọc, hiểu và tuân thủ các quy định trong tài liệu này.

\vfill

\begin{center}
    \textbf{ĐẠI DIỆN BAN QUẢN LÝ ĐÀO TẠO AILMS}\\[1.5cm]
    \textit{Ký và ghi rõ họ tên}
\end{center}

\end{document}